\section{Capítulo~\ref{sec:dendrites}. Número de dendritas}

La estructura de las clases y el número de entradas están medidos anatómicamente y con registros eléctricos; la estimación~\eqref{eq:dendriteload} es física elemental del condensador, y la explicación computacional del número de entradas se apoya en teoría y modelos.

{\footnotesize
\setlength{\tabcolsep}{3pt}\renewcommand{\arraystretch}{1.15}
\begin{longtable}{>{\raggedright}p{0.5cm}>{\raggedright}p{4.0cm}>{\raggedright}p{5.6cm}>{\raggedright}p{2.3cm}>{\raggedright\arraybackslash}p{1.7cm}}
\toprule
N.º & Afirmación & Experimento y qué se midió & Fuente & Estado \\
\midrule\endfirsthead
\toprule
N.º & Afirmación & Experimento y qué se midió & Fuente & Estado \\
\midrule\endhead
1 & Capacidad específica de la membrana $c_m\approx1$~\textmu F/cm$^2$ & Se extrajo membrana del cuerpo de la neurona con una pipeta, se aplicó un escalón de voltaje y, por la caída de la corriente capacitiva, se halló la capacidad total, que luego se dividió por el área: $0.9$~\textmu F/cm$^2$ en tres clases de neuronas & \cite{Gentet2000} & medido \\
2 & Tiempo de carga $t\approx c_mA\Delta V/I$~\eqref{eq:dendriteload}: una neurona grande necesita decenas de entradas simultáneas; a una pequeña le basta una sinapsis grande & Estimación según la ley de carga del condensador; los números ($20$ y $300$~pF, $0.1$ y varios nA) son órdenes de magnitud & deducción del capítulo & teoría \\
3 & Una sola sinapsis enorme da una corriente de varios nanoamperios y lleva de inmediato a una neurona pequeña al umbral & Registro simultáneo de la terminal y de la neurona destinataria en una rodaja: corriente de una sinapsis de $2$--$13$~nA; cada espiga de entrada provoca una respuesta supraumbral, con un retardo de unos $0.4$~ms a $36^\circ$C; la destinataria emite \nt{GLYC} & \cite{Borst1995,BorstSoria2012} & medido \\
4 & Cero dendritas: en las neuronas sensoriales del tacto y del dolor, la espiga va del sensor a la médula espinal sin pasar por el cuerpo celular & La estructura (un axón que se divide en dos junto al cuerpo) está establecida anatómicamente. Que la conducción no requiera que el cuerpo sea excitable se mostró con un modelo validado de esta neurona: sin cuerpo excitable, la espiga atraviesa la bifurcación con la misma fiabilidad & \cite{AmirDevor2003} & teoría \\
5 & Una dendrita: la neurona olfativa lleva un solo tipo de receptor y transmite una sola línea de entrada & Revisión de experimentos con genes de receptores: cada neurona olfativa expresa un único gen de receptor de una gran familia & \cite{Mombaerts2004} & medido \\
6 & Dos dendritas: en los detectores de la dirección del sonido, las entradas de los oídos izquierdo y derecho se posan en dendritas distintas & Anatomía y registros en mamíferos, reunidos en una revisión: las entradas de los dos oídos se reparten en dos dendritas opuestas & \cite{Grothe2010} & medido \\
7 & Repartir las entradas en dos dendritas hace el AND más estricto & Mostrado en un modelo: la saturación~\eqref{eq:shunt} en una dendrita impide que las entradas de un solo lado lleguen al umbral, y la detección de coincidencias mejora. No hay ningún experimento directo en el que se cambie la distribución de las entradas & \cite{AgmonSnir1998} & teoría \\
8 & Unos $7\cdot10^{10}$ pequeñas neuronas \nt{GLUT} del circuito de aprendizaje motor, con $\sim4$ entradas cada una & Recuento de células por fraccionamiento isotrópico: en esta parte del cerebro humano hay unos $69$ mil millones de neuronas. Registro de una de estas células en una rata viva: ante un toque llegan ráfagas de corrientes discretas desde pocas entradas, y la célula dispara cuando se suman & \cite{Azevedo2009,Chadderton2004} & medido \\
9 & Un elemento de umbral con $n$ entradas separa como máximo $P_{\max}\approx2n$ patrones aleatorios~\eqref{eq:cover}; para la neurona de salida del circuito de aprendizaje motor, la cota es $\sim2\cdot10^5$, y la estimación realista, $\sim4\cdot10^4$ asociaciones & La cota es un resultado clásico de la geometría combinatoria. La estimación realista es la teoría de capacidad con pesos solo positivos y margen frente al ruido, aplicada al número medido de entradas de esta neurona & \cite{Cover1965,Brunel2004} & teoría \\
10 & Los mezcladores con pocas entradas sirven para expandir la entrada; “cuatro” está cerca del óptimo & Teoría: la dimensión de la representación que da la capa de mezcladores es máxima aproximadamente con el número de entradas que se observa anatómicamente; la hipótesis original de la expansión está en trabajos clásicos & \cite{Marr1969,Albus1971,LitwinKumar2017} & hipótesis \\
11 & Árboles grandes: $10^4$--$10^5$ entradas & Recuento de sinapsis en micrografías electrónicas: $\approx1.7\cdot10^5$ sinapsis en una sola neurona \nt{GABA} de salida del circuito de aprendizaje motor en la rata; unas $30\,000$ entradas excitadoras y $1\,700$ inhibidoras en una neurona \nt{GLUT} del hipocampo & \cite{NapperHarvey1988,Megias2001} & medido \\
12 & Las ramas de un árbol grande son subcircuitos independientes con sus propios umbrales; la neurona es una pequeña red de varias capas & Estimulación puntual en dos lugares de dendritas finas: las entradas en una misma rama se suman de forma sigmoidea, y en ramas distintas, de forma lineal. En las dendritas de neuronas de la corteza humana se hallaron espigas de calcio graduadas que permiten a una sola célula resolver una tarea linealmente no separable. Modelo: para reproducir una neurona hace falta una red profunda & \cite{Polsky2004,Gidon2020,Beniaguev2021} & medido \\
13 & Dendritas-salida: cada rama de una neurona \nt{GABA} de la retina es un detector de dirección independiente & Registro de calcio con dos fotones en ramas individuales: la señal de la rama es selectiva al movimiento del cuerpo celular hacia la punta, pero el voltaje del cuerpo no & \cite{Euler2002} & medido \\
14 & El número de entradas sigue a la tarea (proposición~\ref{prop:dendrites}) & La estructura de las clases está medida (filas 3--13); que el número de entradas se haya elegido por velocidad o por clasificación solo lo muestran la teoría y los modelos para clases concretas & \cite{LitwinKumar2017,AgmonSnir1998} & hipótesis \\
\bottomrule
\end{longtable}}
